\documentclass[10pt, a4paper]{article}
\usepackage{preamble}

\newcommand\independent{\protect\mathpalette{\protect\independenT}{\perp}}
\def\independenT#1#2{\mathrel{\rlap{$#1#2$}\mkern2mu{#1#2}}}

\title{Statistical Inference II}
\author{Luke Phillips}
\date{August 2025}

\begin{document}

\maketitle

\newpage

\tableofcontents

\newpage

\section{Introduction}

\subsection{Statistical Inference}

\begin{definition}[The general statistical problem]
    We observe the data values $X_1, \dotsc, X_n$ from a distribution $F$.
    We then want to estimate
    (or learn)
    $F$ itself,
    or some characteristics of $F$
    (e.g. its mean).
\end{definition}

\subsubsection{Models and parameters}
A \textbf{statistical model}
(also known as a family of distributions)
is a set of distributions
(or densities)
$\mathcal{M}$.

A \textbf{parametric model} is any model $\mathcal{M}$ that can be parameterised by a finite number of parameters.

In general,
a parametric model for the data would take the form:
\[
\mathcal{M} = \{f(x; \theta)\,:\,\theta \in \Omega\}
\]
where $\theta$ is an unknown parameter
(or vector of parameters)
that can take values in the parameter space $\Omega$.
If $\theta$ is a vector,
we are only interested in one component of $\theta$,
we call the remaining parameters nuisance parameters.

A nonparametric model is a set $\mathcal{M}$ that cannot be parameterised by a finite number of parameters.

\subsubsection{Fundamental Inference Problems}

\textbf{Point Estimation}:

Point estimation refers to providing a single "best guess" for some quantity of interest.
This is often a parameter $\theta$ in a parametric model,
but can also be the entire distribution $F$,
a prediction for a future value $Y$,
or any combination or function of the above.

\textbf{Uncertainty Estimation}:

We need to understand the accuracy of the estimates we make,
we can do this by estimating attributes of our potential error
(such as means and variances),
or producing intervals of plausible values consistent with our estimate.

\textbf{Hypothesis Tests}:

In hypothesis testing,
we start with some default theory about the model - the null hypothesis,
and we ask if the data provides sufficient evidence to reject the theory.

\newpage

\section{Distribution Theory}

\subsection{Multivariate probability}

\begin{definition}[Random vector]
    If $X_1, \dotsc, X_p$ is a collection of real-valued random variables,
    then
    \[
    \mathbf{X} = \begin{bmatrix}
        X_1 \\ \vdots \\ X_p
    \end{bmatrix} = (X_1, \dotsc, X_p) ^ T,
    \]
    is a $p$-dimensional random vector in $\R ^ p$,
    and $\mathbf{x} = (x_1, \dotsc, x_p) ^ T$ is a particular realisation of $\mathbf{X}$.
\end{definition}

\begin{definition}[PDF of a random vector]
    The pdf of a random vector $\mathbf{X} = (X_1, \dotsc, X_p) ^ T$ is the joint probability density function $f(x_1, \dotsc, x_p)$ from $\R ^ p$ to $\R_{0}^{+} = [0, \infty)$ such that:
    
    \begin{enumerate}[label = (\roman*)]
        \item $f(\mathbf{x}) \geq 0$ for all $\mbf{x} \in \R ^ p$

        \item $\displaystyle\int_{\R ^ p}f(\mbf{x})\,\mbf{dx} = \int_{-\infty}^{\infty} \dotsi\int_{-\infty}^{\infty}f(x_1, \dotsc, x_p)\,\mbf{d}x_1, \dotsc, \mbf{d}x_p = 1$

        \item For any subset $\mathcal{A}$ of $\R ^ p\,:\,\P[\mbf{X} \in \mathcal{A}] = \int_{\mathcal{A}}f(\mbf{x})\,\mbf{dx}$.
    \end{enumerate}
\end{definition}

\begin{definition}[CDF of a random vector]
    Similarly,
    the joint cdf is defined as
    \[
    F(\mbf{x}) = F(x_1', \dotsc, x_p') = \P[X_1 \leq x_1', \dotsc, X_p \leq x_p'] = \int_{-\infty}^{x_1'}\dotsi\int_{-\infty}^{x_p'}f(x_1, \dotsc, x_p)\,\mbf{d}x_p, \dotsc, \mbf{d}x_1,
    \]
    and we can recover the pdf by
    (partial)
    differentiation
    \[
    f(\mbf{x}) = \frac{\partial ^ p}{\partial x_1\dotsi\partial x_p}F(\mbf{x}).
    \]
\end{definition}

\subsubsection{Key operations and relations}

\begin{definition}[Marginal pdf]
    A subcollection $(X_1, \dotsc, X_q)$ of $(X_1, \dotsc, X_p)$,
    where $q < p$,
    has a marginal pdf defined as
    \[
    f(x_1, \dotsc, x_q) = \int_{-\infty}^{\infty}\dotsi\int_{-\infty}^{\infty}f(x_1, \dotsc, x_p)\,\mbf{d}x_{q + 1}\dotsc\mbf{d}x_p.
    \]
\end{definition}

\begin{definition}[Conditional pdf]
    If we know the values $X_{q + 1} = x_{q + 1}, \dotsc, X_p = x_p$,
    but not the values of $X_1, \dotsc, X_q$,
    then the $X_1, \dotsc, X_q$ are still random variables and have a conditional pdf:
    \[
    f(x_1, \dotsc, x_q\mid x_{q + 1}, \dotsc, x_p) = \frac{f(x_1, \dotsc, x_q, x_{q + 1}, \dotsc, x_p)}{f(x_{q + 1}, \dotsc, x_p)}.
    \]
\end{definition}

\begin{definition}[Independence]
    A collection of random variables $(X_1, \dotsc, X_p)$ is said to be
    (mutually)
    independent if and only if their joint pdf factorises as
    \[
    f(x_1, x_2, \dotsc, x_p) = f_1(x_1)f_2(x_2)\dotsc f_p(x_p),
    \]
    for all $(x_1, x_2, \dotsc, x_p) \in \R ^ p$.
    We sometimes write this as $X_1 \independent X_2 \dotsi \independent X_p$,
    where $\independent$ means 'independent of'.

    An equivalent statement,
    is that the conditional densities are equal to the marginal densities:
    \[
    f(x_i\mid x_1, \dotsc, x_{i - 1}, x_{i + 1}, \dotsc, x_p) = f(x_i),
    \]
    for all $i = 1, \dotsc, p$.
\end{definition}

\begin{definition}[Random sample,
IID]
    The random variables $X_1, \dotsc, X_p$ are called independent and identically distributed
    (iid)
    random variables with pdf or pmf $f(x)$ if they are all independent of each other,
    and the marginal pdf or pmf of each $X_i$ is the same function $f(x)$.
    This means that the joint pdf simplifies to
    \[
    f(x_1, x_2, \dotsc, x_p) = f(x_1)f(x_2)\dotsc f(x_p) = \prod_{i = 1}^{p}f(x_i).
    \]
    Alternatively,
    $X_1, \dotsc, X_p$ are called a random sample of size $n$ from the distribution
    (or population)
    $f(x)$.
\end{definition}

\subsection{Multivariate expectation and variance}

\begin{definition}[Expectation of a functional]
    Let $g(\mbf{X})\,:\, \R ^ p \to \R$ be a function of the random vector $\mbf{X}$.
    Then $g(\mbf{X})$ is a random variable and its expectation is
    \[
    \E\left(g(\mbf{X})\right) = \int_{-\infty}^{\infty}\dotsi\int_{-\infty}^{\infty}g(x_1, \dotsc, x_p)f(x_1, \dotsc, x_p)\,\mbf{d}x_1\dotsc\mbf{d}x_p = \int_{\R ^ p}g(\mbf{x})f(\mbf{x})\,\mbf{dx}.
    \]
\end{definition}

\begin{definition}[Mean vector]
    The expectation $\E(\mbf{X})$ of a random vector $\mbf{X} = (X_1, \dotsc, X_p)$ is defined as:
    \[
    \E(\mbf{X}) = \int\mbf{x}f(\mbf{x})\,\mbf{dx} = \begin{bmatrix}
        \int_{-\infty}^{\infty}x_1f(\mbf{x})\,\mbf{dx} \\
        \vdots \\
        \int_{-\infty}^{\infty}x_pf(\mbf{x})\,\mbf{dx}
    \end{bmatrix}
    =
    \begin{bmatrix}
        \E(\mbf{X}_1) \\
        \vdots \\
        \E(\mbf{X}_p)
    \end{bmatrix}
    = \mbf{\mu},
    \]
    which is simply the $p$-dimensional vector whose components are the expectations of the individual components of $\mbf{X}$ and is known as the mean vector,
    $\mbf{\mu}$,
    of $\mbf{X}$.
\end{definition}

\begin{definition}[Variance matrix]
    The variance matrix of the random vector $\mbf{X}$ is defined to be the $p \times p$ matrix such that,
    for $i, j = 1, \dotsc p$,
    the $(i, j)$ element of the matrix is $\Cov(X_i, X_j)$:
    \[
    \Var(\mbf{X}) = \E\left((\mbf{X} - \mbf{\mu})(\mbf{X} - \mbf{\mu}) ^ T\right) = \E(\mbf{X}\mbf{X} ^ T) - \mbf{\mu}\mbf{\mu} ^ T = \begin{bmatrix}
        \sum_{11} & \dotsi & \sum_{1p} \\
        \vdots & \ddots & \vdots \\
        \sum_{p1} & \dotsi & \sum_{pp}
    \end{bmatrix} = \mbf{\Sigma}
    \]
    where the components of the matrix comprise the individual variances and covariances of the random variables
    \begin{align*}
        \Sigma_{jj} &= \E(X_j ^ 2) - \E(X_j) ^ 2 = \Var(X_j) = \sigma_j ^ 2 \\
        \Sigma_{ij} &= \E(X_iX_j) - \E(X_i)\E(X_j) = \Cov(X_i, X_j) = \Cov(X_j, X_i) = \sigma_{ji}.
    \end{align*}
\end{definition}

\subsubsection{Properties of expectation and variance}
For shorthand let $\mbf{X} \sim [\mbf{\mu}, \mbf{\Sigma}]$ denote the random vector $\mbf{X}$ which has $\E(\mbf{X}) = \mbf{\mu}$ and variance $\Var(\mbf{X}) = \mbf{\Sigma}$.

\begin{proposition}
    Let $\mbf{X}$ be a random $p$-vector $\mbf{X} \sim [\mbf{\mu}, \mbf{\Sigma}]$.
    Then for constant $p \times p$ matrices $\mbf{A}, \mbf{C}$ and constant $p$-vectors $\mbf{b}, \mbf{d}$:

    \begin{enumerate}[label = (\roman*)]
        \item $\E(\mbf{A}\mbf{X} + \mbf{b}) = \mbf{A}\E(\mbf{X}) + \mbf{b}$,
        
        \item $\Var(\mbf{A}\mbf{X} + \mbf{b}) = \mbf{A}\Var(\mbf{X})\mbf{A} ^ T$,
        
        \item $\Cov(\mbf{A}\mbf{X} + \mbf{b}, \mbf{C}\mbf{X} + \mbf{d}) = \mbf{A}\Var(\mbf{X})\mbf{C} ^ T$.
    \end{enumerate}

    \begin{proof}
        \begin{enumerate}[label = (\roman*)]
            \item
            \begin{align*}
                \E(\mbf{A}\mbf{X} + \mbf{b}) &= \int(\mbf{Ax} + \mbf{b})f(\mbf{x})\,\mbf{dx} + \mbf{b} = \int\mbf{Ax}f(\mbf{x}) + \mbf{b}f(\mbf{x})\,\mbf{dx} = \int\mbf{Ax}f(\mbf{x})\,\mbf{dx} + \int\mbf{b}f(\mbf{x})\,\mbf{dx} \\
                &= \mbf{A}\int\mbf{x}f(\mbf{x})\,\mbf{dx} + \mbf{b} \\
                &= \mbf{A}\E(\mbf{X}) + \mbf{b}.
            \end{align*}

            \item
            Define $\mbf{Y} = \mbf{AX} + \mbf{b}$,
            then by (i) $\mbf{\mu}_Y = \mbf{A}\mbf{\mu} + \mbf{b}$.

            \begin{align*}
                \Var(\mbf{AX} + \mbf{b}) &= \E\left((\mbf{Y} - \mbf{\mu}_Y)(\mbf{Y} - \mbf{\mu}_Y) ^ T\right) = \E\left((\mbf{AX} + \mbf{b} - \mbf{A\mu} - \mbf{b})(\mbf{AX} + \mbf{b} - \mbf{A\mu} - \mbf{b}) ^ T\right) \\
                &= \E\left(\mbf{A}(\mbf{X} - \mbf{\mu})(\mbf{A}(\mbf{X} - \mbf{\mu})) ^ T\right)
                \intertext{Now,
                note that $(\mbf{AB}) ^ T = \mbf{B} ^ T\mbf{A} ^ T$ so:}
                &= \E\left(\mbf{A}(\mbf{X} - \mbf{\mu})(\mbf{X} - \mbf{\mu}) ^ T\mbf{A} ^ T\right)
                \intertext{Then by property (i) we can pull out the constant $\mbf{A}$ on either side to get:}
                &= \mbf{A}\E\left((\mbf{X} - \mbf{\mu})(\mbf{X} - \mbf{\mu}) ^ T\right)\mbf{A} ^ T \\
                &= \mbf{A}\Var(\mbf{X})\mbf{A} ^ T.
            \end{align*}
        \end{enumerate}
    \end{proof}
\end{proposition}

\begin{proposition}
    The variance matrix $\mbf{\Sigma}$ is symmetric and positive semi-definite.
\end{proposition}

\subsection{A change of variables}

\subsubsection{Univariate Transforms}
To look at the distribution and density of $Y$ we can use the following methods:
\begin{enumerate}[label = (\Alph*)]
    \item The cumulative distribution function method:
    find and differentiate the cdf of the transformed random variable to recover the pdf.

    \item The change of variables theorem:
    a general result that is easier to apply when the transformation function $y(\cdot)$ obeys certain properties.
\end{enumerate}

\textbf{The cumulative distribution function method}

This method has two simple steps and works for all general problems where we can find $\P(Y \leq y)$:
\begin{enumerate}[label = \arabic*.]
    \item Find $F_Y(y) = \P(Y \leq y)$ by direct probability calculations using $f_X(x)$ and $y(\cdot)$.

    \item Differentiate $F_Y(y)$ with respect to $y$:
    $f_Y(y) = \frac{d}{dy}F_Y(y)$.
\end{enumerate}

\textbf{The $\chi ^ 2$ distribution}
\begin{example}
    If $Z \sim \mathcal{N}(0, 1)$,
    then $Y = Z ^ 2$ is said to have a chi-squared distribution with one degree of freedom,
    $Y \sim \chi_1 ^ 2$.
    Find the pdf of $Y = Z ^ 2$.

    \begin{solution}
        We start by finding the cdf of $Y$:
        \begin{align*}
            F_Y(y) &= \P(Y \leq y) = \P(Z ^ 2< y) \\
            &= \P(-\sqrt{y} \leq Z \leq \sqrt{y}) \\
            &= \Phi(\sqrt{y}) - \Phi(-\sqrt{y}),
        \end{align*}
        where $\Phi(\cdot)$ is the cdf of the standard normal distribution.
        Note $Z$ can be any real number,
        the transformation $Y = Z ^ 2$ is not one-to-one.
        Thus the probability in the second line is more complicated than the single equality that we might naively expect.
        We then find the pdf by differentiation:
        \begin{align*}
            f_Y(y) &= \frac{\mbf{d}}{\mbf{d}y}F_Y(y) \\
            &= \frac{1}{2}y ^ {-\frac{1}{2}}\phi(\sqrt{y}) + \frac{1}{2}y ^ {-\frac{1}{2}\phi(-\sqrt{y})} \\
            &= y ^ {-\frac{1}{2}\phi(\sqrt{y})} \\
            &= \frac{y ^ {-\frac{1}{2}}}{\sqrt{2\pi}}e ^ {-\frac{1}{2}y},
        \end{align*}
        where $\phi(\cdot)$ is the pdf of the standard normal distribution and we have used the facts that $\phi(z) = \frac{\mbf{d}}{\mbf{d}z}\Phi(z)$,
        and $\phi(z)$ is symmetric about $z = 0$.
        Thus $f_Y(y)$ is the pdf of the $\chi_1 ^ 2$ distribution.
    \end{solution}
\end{example}

\textbf{Properties:}

\begin{enumerate}[label = \arabic*.]
    \item If $U \sim \chi_1 ^ 2$ then $\E(U) = 1$ and $\Var(U) = 2$.
    
    \item If $U_1, \dotsc, U_k$ are all iid $\chi_1 ^ 2$ and $V = \sum_{i = 1}^{k}U_i$,
    then $V \sim \chi_k ^ 2$.

    \item If $V \sim \chi_k ^ 2$ then $\E(V) = k$ and $\Var(V) = 2k$.
\end{enumerate}

\textbf{Univariate transformation theorem}

If the function $y(X)$ is monotonic,
invertible and differentiable over the support,
$\chi$,
of $X$,
then we can find a general method to determine the pdf of any such transformation of $X$.

\begin{theorem}[Distribution of a univariate transformation]
    Let $X$ be a continuous random variable with probability density function $f_X(x)$ with support $\chi = [a, b]$.
    Let $Y = y(X)$ where $y(\cdot)$ is a monotonic,
    invertible and differentiable function over $\chi$ with inverse $X = x(Y) = y ^ {-1}(Y)$.
    Then $Y$ has a probability density function $f_Y(y)$ given by
    \[
    f_Y(y) = f_X(x(y))\left|\frac{\mbf{d}x(y)}{\mbf{d}y}\right|,
    \]
    for $y$ in $\mathcal{Y} = [y(a), y(b)]$.

    \begin{proof}
        Using the cdf method to prove this.
        If $Y = y(X)$,
        then the cdf of $Y$ is:
        \[
        F_Y(y) = \P(Y \leq y) = \P(y(X) \leq y).
        \]
        We are given that $y(\cdot)$ is monotonic,
        which leads to two cases:
        $y(\cdot)$ is monotonically increasing,
        $y(\cdot)$ is monotonically decreasing.

        Assume $y(\cdot)$ is increasing,
        then
        \[
        \P(y(X) \leq y) = \P(X \leq y ^ {-1}(y)) = F_X(x(y)),
        \]
        by the definition of the cdf of $X$.
        Now,
        get the pdf by differentiating
        \[
        f_Y(y) = \frac{d}{dy}F_y(y) = \frac{d}{dy}F_X(x(y)) = f_X(x(y))\frac{dx(y)}{dy},
        \]
        by the chain rule.

        For $y(\cdot)$ decreasing,
        we find
        \[
        F_Y(y) = \P(X \geq x(y)) = 1 - F_X(x(y))
        \]
        giving the pdf as
        \[
        f_Y(y) = -f_X(x(y))\frac{dx(y)}{dy},
        \]
        since $x(y)$ is a decreasing function we get $\frac{d}{dy}x(y) < 0$,
        which cancels the minus sign.
        So we can combine the two equations to get the result
        \[
        f_Y(y) = f_X(x(y))\left|\frac{\mbf{d}x(y)}{\mbf{d}y}\right|.
        \]
    \end{proof}
\end{theorem}

\subsubsection{Multivariate Transforms}

\begin{theorem}[Distribution of a multivariate transformation]
    Suppose we have a random $p$-vector $\mbf{X} = (X_1, \dotsc, X_p) ^ T$ with joint continuous pdf $f_{\mbf{X}}(x_1, \dotsc, x_p)$.
    Suppose we create a transformed random variable $\mbf{Y} = (Y_1, \dotsc, Y_p) ^ T$ where each $Y_i$ is a transformation of the $\mbf{X}\,:\, Y_1 = y_1(\mbf{X}), \dotsc, Y_p = y_p(\mbf{X})$.
    So,
    $\mbf{Y} = y(\mbf{X})$,
    where we assume that $y(\cdot)$ is an invertible differentiable transformation with inverse
    \[
    \mbf{X} = y ^ {-1}(\mbf{Y}) = x(\mbf{Y}) = (x_1(\mbf{Y}), \dotsc, x_p(\mbf{Y})).
    \]
    Let the support of the pdf of $X$ be $\chi \subseteq \R ^ p$ then as $X$ varies over $\chi$,
    $Y$ will vary over some set $\mathcal{Y}$ thus $\mathcal{Y}$ is the image of $\chi$.
    Finally,
    let us assume that all of the partial derivatives $\frac{\partial x_i}{\partial y_j}$ exist for all $i, j$ and $y \in \mathcal{Y}$.

    Given $\mbf{X}, \mbf{Y}, y(\mbf{X}), x(\mbf{Y})$ as above,
    then the joint density for $\mbf{Y}$ is given by:
    \[
    f_{\mbf{Y}}(\mbf{y}) = \begin{cases}
        f_{\mbf{X}}(x(\mbf{y}))|\det(\mbf{J})|, & y \in \mathcal{Y} \\
        0, & \text{otherwise}
    \end{cases}
    \]
    where $\mbf{J}$ is the Jacobian matrix of the transformation:
    \[
    \mbf{J} = \begin{bmatrix}
        \frac{\partial x_1}{\partial y_1} & \dotsi & \frac{\partial x_1}{\partial y_p} \\
        \vdots & \ddots & \vdots \\
        \frac{\partial x_p}{\partial y_1} & \dotsi & \frac{\partial x_p}{\partial y_p}
    \end{bmatrix}
    = \left[\frac{\partial x_i}{\partial y_j}\right]_{ij} = \left[\frac{\partial\mbf{x}}{\partial\mbf{y}}\right].
    \]
    
\end{theorem}










\end{document}